\section{Proofs and Additional Experimental Details}

\subsection{Proof of \cref{prop:no-overlap}}

Because \(P_X(A)=0\), the labeled source sample contains no points from \(A\). The target covariates reveal the mass \(Q_X(A)=\pi\), but they reveal no labels in \(A\). Construct two worlds that agree on \(P_X\), \(Q_X\), and \(P(Y\given X)\) outside \(A\). Inside \(A\), set
\[
\EE_{Q_0}[Y-s(X)\given X=x]=0
\]
in world 0 and
\[
\EE_{Q_1}[Y-s(X)\given X=x]=\Delta
\]
in world 1. Both worlds satisfy covariate shift within their own joint laws because the source conditional label law can be defined to match the target conditional law on the support of \(P_X\), and \(P_X(A)=0\). The observable distribution of source labels and target covariates is identical, while the target cell residuals differ by \(\Delta\). Therefore no procedure based on the observed data can uniformly distinguish the two cases.

\subsection{Proof of \cref{thm:main-certificate}}

Fix the source and target covariates and any weight-fitting randomness. Under \cref{ass:label-independent-weights}, the labels are independent and the audit weights \(v_i\) are fixed. For a cell \(c\), write
\[
\alpha_{i,c}
=
\frac{v_iA_c(X_i)}{\sum_jv_jA_c(X_j)} .
\]
Then
\[
\widehat r_c(v)=\sum_i\alpha_{i,c}Z_i,
\qquad
\bar r_c(v)=\sum_i\alpha_{i,c}m_i,
\qquad
\sum_i\alpha_{i,c}^2=\frac{1}{\widehat{\ess}_c(v)}.
\]
Conditional on \(X_i\), the Bernoulli residual \(Z_i=Y_i-s(X_i)\) takes values in an interval of length one. Hoeffding's inequality for weighted sums gives
\[
\PP\left(
\left|\widehat r_c(v)-\bar r_c(v)\right|>t
\right)
\le
2\exp\{-2t^2\widehat{\ess}_c(v)\}.
\]
Taking
\[
t=\sqrt{\frac{\log(2M/\delta_1)}{2\widehat{\ess}_c(v)}}
\]
and applying a union bound over \(M\) cells proves the first part.

For learned weights, let
\[
\widehat S_c=\sum_i\widehat w_iA_c(X_i),
\qquad
S_c=\sum_iw_iA_c(X_i),
\qquad
N_c=\sum_iw_iA_c(X_i)m_i .
\]
Because \(|m_i|\le1\),
\[
\left|
\sum_i(\widehat w_i-w_i)A_c(X_i)m_i
\right|\le\Delta_c,
\qquad
|\widehat S_c-S_c|\le\Delta_c,
\qquad
\widehat S_c\ge S_c-\Delta_c.
\]
The elementary ratio inequality
\[
\left|
\frac{\widehat N}{\widehat S}-\frac{N}{S}
\right|
\le
\frac{|\widehat N-N|}{\widehat S}
+
\frac{|N|\,|\widehat S-S|}{S\widehat S}
\]
with \(|N|\le S\) yields
\[
\left|\bar r_c(\widehat w)-\bar r_c(w)\right|
\le
\frac{2\Delta_c}{S_c-\Delta_c}
\le \rho_c .
\]
Adding this deterministic perturbation to the label-noise event proves the learned-weight statement.

For the population term, assume \(0\le w(X)\le W\). The variables
\[
w(X_i)A_c(X_i)m_i
\quad\text{and}\quad
w(X_i)A_c(X_i)
\]
both lie in intervals of length at most \(2W\) and \(W\), respectively, and the first has absolute value at most \(W\). A finite-family Hoeffding bound gives, with probability at least \(1-\delta_2\),
\[
\left|
\frac1n\sum_iw(X_i)A_c(X_i)m_i
-
\EE_P[w(X)A_c(X)m(X)]
\right|
\le \epsilon_n(\delta_2)
\]
and
\[
\left|
\frac1n\sum_iw(X_i)A_c(X_i)-\mu_c^Q
\right|
\le \epsilon_n(\delta_2)
\]
simultaneously over cells, after adjusting constants. On this event, the same ratio perturbation argument gives
\[
|\bar r_c(w)-r_c^Q|
\le
\frac{2\epsilon_n(\delta_2)}{\mu_c^Q-\epsilon_n(\delta_2)}
=
\gamma_c(\delta_2).
\]
Combining the label, learned-weight, and population events gives probability at least \(1-\delta_1-\delta_2\).

\subsection{Proof of \cref{cor:certified-alarm}}

On the event from \cref{thm:main-certificate},
\[
|\widehat r_c(\widehat w)-r_c^Q|\le B_c
\]
simultaneously over \(c\). If \(|r_c^Q|\le\tau\) for every cell, then
\[
|\widehat r_c(\widehat w)|-B_c\le |r_c^Q|\le\tau
\]
for every cell. Therefore no certified alarm occurs on that event. The failure probability is at most \(\delta_1+\delta_2\).

\subsection{Proof of \cref{prop:global-misleading}}

Take one active point, \(A_1=1\), and \(N-1\) inactive points, \(A_2=\cdots=A_N=0\). Assign
\[
w_1=N\pi,
\qquad
w_2=\cdots=w_N=\frac{N(1-\pi)}{N-1}.
\]
The weighted mass of the active cell is \(\pi\), and its local ESS is \(1\). The global ESS is
\[
\frac{(N\pi+N(1-\pi))^2}
{(N\pi)^2+(N-1)\{N(1-\pi)/(N-1)\}^2}
=
\frac{1}{\pi^2+(1-\pi)^2/(N-1)}.
\]

\subsection{Additional Experimental Details}

The synthetic certification experiment uses oracle density-ratio weights for the Gaussian shift benchmark. The null setting is calibrated by construction. The planted setting adds a residual shift of \(0.60\) on \(x_0>0.5\), clipped to keep probabilities in \([10^{-3},1-10^{-3}]\). The four rules are the naive worst-cell threshold, a global ESS gate at \(50\), a local ESS gate at \(20\), and the simultaneous local certificate.

A separate local-overlap diagnostic run, reported in earlier drafts and preserved in the repository, shows that lower-tail local ESS collapses much faster than global ESS. At shift \(0.8\), the median global ESS is about \(95\), the lower-tail local ESS is about \(1.4\), and the false-alarm rate above \(0.05\) is about \(80\%\).

The learned-weight experiment compares plugin logistic density-ratio weights and source-side cross-fitted logistic weights. Because the oracle ratio is known in this benchmark, we compute the nuisance term
\[
\frac{2\Delta_c}{S_c-\Delta_c}
\]
whenever \(S_c>\Delta_c\), and set it to \(1\) otherwise. This deliberately conservative oracle correction is used only to show how local density-ratio error can dominate a certificate.

\subsection{Constructing \texorpdfstring{\(\widehat\rho_c\)}{rho-c} in Practice}

The population quantity \(\rho_c\) depends on the unknown oracle ratio. A deployable audit should therefore distinguish three levels of evidence.

\textbf{Model-based certificate.} If the density-ratio procedure comes with a high-probability local log-error bound
\[
\left|\log \widehat w_i-\log w(X_i)\right|\le\kappa_c
\qquad\text{for all }i\text{ with }A_c(X_i)=1,
\]
then the main text gives the plug-in envelope
\[
\widehat\rho_c
=
\frac{2(e^{\kappa_c}-1)}{2-e^{\kappa_c}},
\qquad
\kappa_c<\log 2.
\]
Such a bound might come from a correctly specified parametric ratio model, a conservative uncertainty band for the domain classifier's logit, or an externally validated density-ratio model. If \(\kappa_c\ge\log2\), the bound is too weak to certify the cell.

\textbf{Sensitivity envelope.} More commonly, the audit team has no theorem-level guarantee for the ratio model. The practical alternative is to construct a candidate set \(\mathcal W\) of plausible cross-fitted ratios. Useful perturbations include:
\begin{itemize}
  \item varying the domain-classifier class or regularization strength;
  \item varying clipping thresholds for extreme weights;
  \item refitting over bootstrap or subsampled source/target covariates;
  \item adding or removing covariate blocks whose shift mechanism is uncertain;
  \item comparing direct density-ratio methods with probabilistic domain-classifier ratios.
\end{itemize}
For a baseline \(\widehat w\), compute
\[
\widehat\rho_c^{\mathrm{sens}}
=
\max_{\widetilde w\in\mathcal W}
\frac{
2\sum_i|\widehat w_i-\widetilde w_i|A_c(X_i)
}{
\sum_i\widetilde w_iA_c(X_i)
-
\sum_i|\widehat w_i-\widetilde w_i|A_c(X_i)
},
\]
omitting candidates with nonpositive denominators and marking the cell unsupported if all denominators fail. This quantity is a rigorous nuisance term only under the additional assumption that \(\mathcal W\) contains the oracle ratio with high probability. Without that assumption, it is still valuable as a transparent stress test: it tells the reader how much the cell conclusion changes across plausible ratio models.

\textbf{Balance diagnostics.} Weight sensitivity should be accompanied by covariate balance checks on held-out source and target covariates. For each decision-driving cell, compare the target cell mass
\[
\widehat\mu_c^Q=\frac1m\sum_{j=1}^m A_c(X'_j)
\]
with the weighted source estimate
\[
\widehat\mu_c^{P,\widehat w}
=
\frac1n\sum_{i=1}^n \widehat w_iA_c(X_i),
\]
and repeat the comparison for important covariate summaries inside the cell. Large imbalance does not directly equal \(\rho_c\), but it is evidence that the ratio model is not transporting the local cell reliably. The appropriate report is then insufficient evidence unless the alarm survives a wider sensitivity envelope.

In all cases, \(\widehat\rho_c\) should be reported only for cells that could affect the decision: the selected worst cell, any certified cells, and a small set of near-threshold cells. Reporting a single global nuisance number hides the same local problem that motivates the paper.

The adaptive-search experiment expands the subgroup class to threshold rules on all six synthetic features and pairwise positive-tail intersections over the first four features, producing roughly \(500\)--\(600\) nonempty cells after filtering. The sample-split protocol uses half of the source audit labels for discovery and half for the final one-cell certificate.

The Adult Census case study uses \texttt{fetch\_openml("adult", version=2)} and removes \texttt{fnlwgt}. The target deployment split is sampled with probability increasing in age, hours worked, and college-level education, with a mild negative tilt for the female subgroup. A logistic income model is trained on \(16000\) source-training examples. The audit uses \(9000\) labeled source-audit examples, \(12000\) target covariates, and cross-fitted logistic domain weights. Target labels are used only for retrospective evaluation.

The ACSIncome case study uses Folktables \citep{Ding2021RetiringAdult}. The source population is California ACS records and the target population is New York ACS records from the 2018 one-year person survey. A logistic income model is trained on \(30000\) California records and audited on \(18000\) held-out California source-audit records against \(18000\) New York target covariates. The subgroup family uses sex, race, age, hours worked, education, and simple intersections. The target labels are held out from the audit and used only for retrospective validation. The learned-weight sensitivity envelope is built from alternative cross-fitted domain-ratio models with different regularization strengths and clipping thresholds.

\subsection{Secondary Results Kept Out of the Main Text}

The main text focuses on certification, but the repository also contains subgroup-degradation examples showing why worst-group auditing is needed in the first place. In representative synthetic scenarios, masked subgroup shift gives a worst-group ECE gap of about \(0.182\), and minority subgroup degradation gives a gap of about \(0.217\). These examples are not used as headline evidence because the paper's contribution is not the existence of hidden subgroup failures; it is the certification problem under covariate shift.

The same reason explains why doubly robust estimators are discussed but not foregrounded. For a residual \(Z=Y-s(X)\), a typical augmented residual-mass estimator has the form
\[
\widehat\theta_c^{\mathrm{DR}}
=
\frac1m\sum_{j=1}^{m}\widehat m(X'_j)A_c(X'_j)
+
\frac1n\sum_{i=1}^{n}\widehat w_iA_c(X_i)
\{Z_i-\widehat m(X_i)\}.
\]
If \(\widehat m\) is accurate, this can reduce nuisance bias and variance. However, the influence-function variance still contains a local weighted residual term,
\[
\EE_P\!\left[
w(X)^2A_c(X)\Var(Z\given X)
\right],
\]
and the no-overlap result remains unchanged. Augmentation can improve constants; it cannot identify labels in an unsupported target cell.

The main rendered artifacts are:
\begin{itemize}
  \item \path{preprint/figures/certified_audit_main.png}
  \item \path{preprint/figures/estimated_weight_and_search.png}
  \item \path{preprint/figures/adult_case_study.png}
  \item \path{preprint/figures/acs_external_validation.png}
  \item \path{preprint/research/results/certified_audit/}
  \item \path{preprint/research/results/estimated_weights/}
  \item \path{preprint/research/results/adaptive_search/}
  \item \path{preprint/research/results/adult_case_study/}
  \item \path{preprint/research/results/acs_external_validation/}
\end{itemize}
